\chapter{Stand der Forschung}
\label{chap:forschungsstand}

\section{Recherchevorgehen}

Die Literaturrecherche wurde in Google Scholar, IEEE Xplore, dem National Bureau of Economic Research sowie Veröffentlichungen der Europäischen Zentralbank durchgeführt. Zusätzlich wurden regulatorische Dokumente der EU und der VAE analysiert. Berücksichtigt wurden Quellen aus dem Zeitraum 2018–2026.

\section{Bisherige Forschung und Methoden}
\label{sec:bisherige-forschung-methoden}

Die relevante Forschung lässt sich in drei Bereiche einteilen.

\textbf{Geldtheoretische Perspektive:} Ammous \hyperref[Ammous2018]{(2018)} zeigt anhand der drei klassischen Geldfunktionen -- Tauschmittel, Recheneinheit und Wertaufbewahrungsmittel --, dass die hohe Volatilität von Kryptowährungen ihre Nutzung als Recheneinheit und Wertaufbewahrungsmittel einschränkt und das nötige Vertrauen als Tauschmittel fehlt. Kryptowährungen werden daher überwiegend als Spekulationsobjekte gehalten und selten als Zahlungsmittel genutzt~\cite{Ammous2018}.

\textbf{Verhaltensökonomie:} Ballis \& Verousis \hyperref[Ballis2022]{(2022)} identifizieren spekulatives Verhalten und Herdeneffekte auf Kryptomärkten. EZB-Daten von Zamora-Pérez \hyperref[Zamora2026]{(2026)} bestätigen dies: Krypto-Assets werden verbreitet gehalten, ihre Nutzung als Zahlungsmittel bleibt jedoch gering~\cite{Ballis2022, Zamora2026}.

\textbf{Regulierung und staatliche Strategien:} Die Forschung zeigt drei typische Ansätze: vollständige Einführung (El Salvador), wirtschaftsorientierte Regulierung (VAE) und Kontrolle ohne Zahlungsfunktion (EU). Alvarez et al. \hyperref[Alvarez2023]{(2023)} zeigen anhand einer Haushaltsbefragung in El Salvador, dass gesetzliche Anerkennung allein keine nachhaltige Nutzung erzeugt~\cite{Alvarez2023}. Arbeiten zu den VAE und der EU belegen, dass Regulierung vor allem Marktstrukturen beeinflusst, nicht aber die alltägliche Zahlungsnutzung~\cite{VARA2025, MiCA2023}.

\subsection{Forschungslücke und Beitrag dieser Arbeit}
\label{subsec:forschungsluecke}

Bisherige Forschung betrachtet Länder meist isoliert oder fokussiert einzelne Aspekte wie Adoption oder Regulierung. Ein systematischer Vergleich unterschiedlicher staatlicher Strategien hinsichtlich ihrer Auswirkungen auf die tatsächliche Zahlungsnutzung fehlt. Diese Arbeit schließt diese Lücke durch einen direkten Vergleich von El Salvador, den VAE und der EU.